% optimization/hybrid_block_structure.tex

\section{Hybrid Representation and Block Structure}

% ============================================================
\subsection{Motivation}
% ============================================================

Purely parametric models provide a low-dimensional and interpretable
representation of alignments, but lack flexibility when fitting
real-world data.

Purely sampled models provide maximum flexibility, but lead to
high-dimensional and often ill-conditioned optimization problems.

A hybrid formulation combines both advantages.

% ============================================================
\subsection{Hybrid Representation}
% ============================================================

\begin{definition}[Hybrid curvature model]
\[
\kappa(s)
=
\kappa_{\mathrm{param}}(s;\mathbf{T})
+
\delta\kappa(s),
\]
where:
\begin{itemize}
\item \(\kappa_{\mathrm{param}}\) is a structured parametric model,
\item \(\delta\kappa(s)\) is a correction term.
\end{itemize}
\end{definition}

\begin{definition}[Hybrid cant model]
\[
\phi(s)
=
\phi_{\mathrm{param}}(s;\mathbf{T}_\phi)
+
\delta\phi(s).
\]
\end{definition}

% ============================================================
\subsection{Discretization of Corrections}
% ============================================================

The correction terms are represented discretely:
\[
\delta\kappa_i = \delta\kappa(s_i),
\qquad
\delta\phi_i = \delta\phi(s_i).
\]

\begin{definition}[Extended parameter vector]
\[
\mathbf{x}
=
(\mathbf{T}, \mathbf{T}_\phi, \delta\boldsymbol{\kappa}, \delta\boldsymbol{\phi}).
\]
\end{definition}

% ============================================================
\subsection{Regularization}
% ============================================================

To avoid overfitting, correction terms are regularized:
\[
J_{\delta}
=
\int_0^L (\delta\kappa'(s))^2 \,\dd s
+
\int_0^L (\delta\phi'(s))^2 \,\dd s.
\]

Thus, the correction captures only smooth, physically meaningful
deviations.

% ============================================================
\subsection{Residual Structure}
% ============================================================

The residual vector consists of:
\begin{itemize}
\item parametric contributions,
\item correction contributions,
\item physical constraints,
\item data fitting terms.
\end{itemize}

% ============================================================
\subsection{Block Structure of Variables}
% ============================================================

The parameter vector decomposes naturally as:
\[
(\mathbf{T}, \delta\boldsymbol{\kappa})
\quad \text{and} \quad
(\mathbf{T}_\phi, \delta\boldsymbol{\phi}).
\]

This yields a two-level structure:
\begin{itemize}
\item coarse scale: parametric variables,
\item fine scale: correction variables.
\end{itemize}

% ============================================================
\subsection{Jacobian Block Structure}
% ============================================================

\begin{definition}[Block Jacobian]
The Jacobian takes the form:
\[
J_r =
\begin{pmatrix}
J_{\kappa\kappa}^{\mathrm{param}} & J_{\kappa\kappa}^{\mathrm{corr}} & 0 & 0 \\
0 & 0 & J_{\phi\phi}^{\mathrm{param}} & J_{\phi\phi}^{\mathrm{corr}} \\
J_{\mathrm{dyn},\kappa} & J_{\mathrm{dyn},\kappa}^{\delta} &
J_{\mathrm{dyn},\phi} & J_{\mathrm{dyn},\phi}^{\delta}
\end{pmatrix}.
\]
\end{definition}

Key properties:
\begin{itemize}
\item parametric and correction terms are weakly coupled,
\item curvature and cant are largely separable,
\item coupling occurs mainly through dynamics.
\end{itemize}

% ============================================================
\subsection{Hierarchical Interpretation}
% ============================================================

The hybrid model introduces a hierarchical structure:
\begin{itemize}
\item parametric layer defines global alignment structure,
\item correction layer refines local deviations.
\end{itemize}

This reflects engineering practice:
\begin{itemize}
\item design defines intent,
\item construction introduces deviations.
\end{itemize}

% ============================================================
\subsection{Computational Strategy}
% ============================================================

\begin{itemize}
\item optimize parametric variables first,
\item introduce correction terms for refinement,
\item optionally alternate between both levels.
\end{itemize}

% ============================================================
\subsection{Connection to Level-of-Detail}
% ============================================================

The hybrid representation naturally supports LOD concepts:
\begin{itemize}
\item coarse LOD: parametric model only,
\item fine LOD: parametric + corrections.
\end{itemize}

% ============================================================
\subsection{Key Insight}
% ============================================================

\begin{proposition}
The hybrid formulation separates structural modelling from data fitting
while preserving a sparse and well-conditioned optimization problem.
\end{proposition}

% ============================================================
\subsection{Connection to AIM}
% ============================================================

\begin{summary}
The combination of hybrid modelling and block structure is a central
feature of the AIM framework.

It enables:
\begin{itemize}
\item interpretable parametric modelling,
\item robust handling of imperfect data,
\item efficient optimization due to sparsity and locality.
\end{itemize}

Thus, it bridges the gap between engineering design and data-driven
reconstruction.
\end{summary}
